% Validation Methodology Section (Revised)
% Corrected methodology: holdout-only, unsupervised threshold, PCA robustness

\subsection{Validation Methodology}
\label{sec:validation_methodology}

We perform comprehensive validation to ensure our model preference heterogeneity findings are statistically robust and not artifacts of methodological choices. All analyses use the held-out test set exclusively ($N=750$), with the PCA model trained on a disjoint set and the threshold determined by purely unsupervised criteria.

\subsubsection{Statistical Significance Testing}

Since reward gaps are derived from discrete pairwise outcomes (win/tie/loss), we use categorical and non-parametric tests as the primary analysis:

\paragraph{Primary Tests (Categorical).}
\begin{itemize}[leftmargin=*,noitemsep,topsep=0pt]
    \item \textbf{$\chi^2$ test of independence:} $p < 0.0001$ (outcome proportions differ significantly between clusters)
    \item \textbf{Cram\'er's $V$:} Reported in figure (categorical effect size)
    \item \textbf{Contingency table:} Win/tie/loss counts by cluster (see figure panel B)
\end{itemize}

\paragraph{Supplementary Tests (Ordinal).}
\begin{itemize}[leftmargin=*,noitemsep,topsep=0pt]
    \item \textbf{Mann-Whitney $U$-test (non-parametric):} $p < 0.0001$ (appropriate for ordinal reward gaps)
    \item \textbf{Cohen's $d$:} $1.53$ with domain-adapted PCA, $0.33$ with generic PCA (approximate---reward gaps are discrete, not continuous)
\end{itemize}

\paragraph{Distribution Diagnostics.}
Shapiro-Wilk normality tests indicate both distributions are non-normal ($p < 0.001$), consistent with the discrete nature of pairwise preference outcomes and justifying the use of non-parametric tests.

\subsubsection{Threshold Validation}

The cluster boundary $\text{PC1} = 0.222$ was determined using a purely unsupervised method: silhouette-optimal threshold search over PC1 values, maximizing geometric cluster quality (silhouette $= 0.499$) without reference to reward labels. This eliminates the methodological concern of reward-based threshold selection bias present in grid-search approaches that incorporate gap separation.

\paragraph{Unsupervised Clustering Convergence.}
Independent unsupervised methods identify consistent structure:
\begin{itemize}[leftmargin=*,noitemsep,topsep=0pt]
    \item Silhouette-optimal: $\text{PC1} = 0.222$ (silhouette $= 0.499$)
    \item $k$-means ($k=2$): boundary $\approx 0.138$
\end{itemize}
Both methods yield $p < 0.0001$.

\paragraph{Statistical Power.}
We conduct a Monte-Carlo power analysis for the $\chi^2$ test given the observed contingency table proportions and sample sizes ($n_{\text{low}} = 606$, $n_{\text{high}} = 144$). With $N = 750$ and the observed outcome proportions, the power to detect the observed effect at $\alpha = 0.05$ exceeds 0.99. For detectable effect sizes at 80\% power, effects as small as Cram\'er's $V \approx 0.10$ are detectable, confirming the sample is adequately powered for this comparison.

\paragraph{Effect Size Stability.}
Since $p$-values alone are uninformative for threshold robustness (with $N=750$ discrete outcomes, most non-trivial splits produce small $p$-values), we additionally report Cram\'er's $V$ across a sweep of candidate thresholds to verify that the effect size---not merely its significance---is stable. The effect size remains meaningful across a range of thresholds, confirming the finding is not an artifact of a particular boundary choice.

\subsubsection{PCA Provenance Sensitivity}

To validate that the finding is not an artifact of domain-adapted PCA training, we repeat the analysis with PCA trained on 100K generic web text samples (C4 corpus). The generic PCA provides an unbiased baseline with no connection to routing:

\begin{table}[h]
\centering
\caption{PCA Provenance Comparison on Held-Out Data ($N=750$)}
\label{tab:pca_provenance}
\begin{tabular}{lccl}
\toprule
\textbf{PCA Training Data} & \textbf{$p$-value} & \textbf{Cohen's $d$} & \textbf{Interpretation} \\
\midrule
C4 web text (100K, generic) & $< 0.0001$ & 0.33 & Small effect (unbiased baseline) \\
Routing prompts (80K, domain-adapted) & $< 0.0001$ & 1.53 & Large effect (amplified by task-relevant features) \\
\bottomrule
\end{tabular}
\end{table}

\noindent \textbf{Key observation:} The \textbf{primary, unbiased effect size} is $d = 0.33$ (small effect, generic PCA), confirming the structure exists independently of PCA provenance. The domain-adapted PCA concentrates routing-relevant variation into PC1, amplifying the effect $4.6\times$ to $d = 1.53$. This amplification is expected---the routing PCA was trained on 80K prompts from the same model-comparison domain---but partly reflects domain overlap between PCA training data and the holdout set, not only genuine signal concentration. The generic result ($d = 0.33$) should be considered the baseline for the effect's true magnitude. Both PCAs are unsupervised (no reward labels).

\subsubsection{High-Dimensional Structure Validation}

\begin{table}[h]
\centering
\caption{Cluster Quality Across Dimensionalities}
\label{tab:multidim_validation}
\begin{tabular}{lccc}
\toprule
\textbf{Space} & \textbf{Dimensions} & \textbf{Silhouette} & \textbf{Separation Ratio} \\
\midrule
2D PCA & 2 & 0.499 & --- \\
32D PCA & 32 & $\sim$0.25 & $>$1.0 \\
384D Embeddings & 384 & 0.057 & 0.81 \\
\bottomrule
\end{tabular}
\end{table}

\noindent \textbf{Important caveat:} The 384D silhouette score of 0.057 indicates weak cluster structure in the full embedding space. The separation ratio of 0.81 ($< 1.0$) means within-cluster distances exceed between-cluster distances. PC1 extracts a moderate task-relevant signal (Spearman $\rho = -0.395$, $\rho^2 = 0.156$, explaining $\sim$16\% of reward gap variance) from a noisy high-dimensional space, not a dominant clustering pattern.

\paragraph{Projection Artifact Consideration.}
Any linear projection of 384D data onto 2D will find \emph{some} direction that correlates with a binary label, raising the concern that the observed structure could be a projection artifact rather than genuine high-dimensional clustering. Two pieces of evidence argue against this interpretation: (1)~the effect persists with generic PCA trained on unrelated C4 web text ($d = 0.33$, $p < 0.0001$), and (2)~the moderate Spearman correlation ($\rho = -0.395$) exceeds what would be expected from random projections. However, this alternative explanation cannot be fully excluded, and the weak 384D structure (silhouette $= 0.057$) means the ``clusters'' should be interpreted as regions of the PC1 axis with different outcome distributions, not as well-separated groups in embedding space.

\subsubsection{Data Quality Analysis}

The held-out dataset ($N=750$) was created using stratified sampling across three dimensions (category, complexity, difficulty) to ensure representative coverage of diverse prompt characteristics (see Appendix~\ref{app:stratification} for complete methodology). The dataset contains zero exact duplicates. Intra-cluster diversity:
\begin{itemize}[leftmargin=*,noitemsep,topsep=0pt]
    \item High PC1 diversity: 0.355 (relatively homogeneous---coherent semantic category)
    \item Low PC1 diversity: 0.953 (highly diverse---broad range of prompt types)
\end{itemize}

\subsubsection{Summary of Validation Approach}

\begin{enumerate}[leftmargin=*,noitemsep,topsep=0pt]
    \item \textbf{Categorical analysis:} $\chi^2$ test on win/tie/loss contingency table ($p < 0.0001$), appropriate for discrete pairwise outcomes
    \item \textbf{Ordinal confirmation:} Mann-Whitney $U$ ($p < 0.0001$); Cohen's $d$ reported as approximate ($d = 0.33$ generic, $1.53$ domain-adapted)
    \item \textbf{Primary effect size:} $d = 0.33$ (generic PCA, unbiased baseline); $d = 1.53$ is the amplified, domain-adapted result
    \item \textbf{Threshold justification:} Purely unsupervised (silhouette-optimal, no reward peeking)
    \item \textbf{PCA sensitivity:} Generic C4 PCA provides unbiased baseline ($d = 0.33$); domain-adapted PCA amplifies to $d = 1.53$ ($4.6\times$), partly reflecting domain overlap
    \item \textbf{High-D caveat:} Weak 384D structure (silhouette $= 0.057$); PC1 captures moderate signal ($\rho^2 = 0.156$); projection artifact alternative discussed
    \item \textbf{Data quality:} Zero duplicates; High PC1 cluster is homogeneous (diversity $= 0.355$); cluster content characterized via TF-IDF analysis
    \item \textbf{Data integrity:} Held-out only ($N=750$); dev set excluded
    \item \textbf{Generalizability:} Cluster proportion (19.2\% High PC1) is specific to the holdout; 1M dataset shows $\sim$5.9\% in corresponding spatial region
\end{enumerate}
